\subsection{Criteria of finiteness for graded rings}

In this no. and the following, all the graduations considered (Algebra, Chapter II, § 11) are assumed to be of type Z. If A (resp. M) is a graded ring (resp. graded module), \(A_{i}\) (resp. \(M_{i}\) ) will denote the set of homogeneous elements of degree i in A (resp. M).

If \(A_{i}=\{0\}\) (resp. \(M_{i}=\{0\}\) ) for i\textless0, A (resp. M) will, to abbreviate, be called a graded ring (resp. module) with positive degrees.

PROPOSITION 1. Let \(\mathbf{A} = \bigoplus_{i\in \mathbb{Z}}\mathbf{A}_i\) be a graded commutative ring with positive degrees, \(\mathfrak{m}\) the graded ideal \(\bigoplus_{i\geqslant 1}\mathbf{A}_i\) and \((x_{\lambda})_{\lambda \in \mathbf{L}}\) a family of homogeneous elements of \(\mathbf{A}\) of degrees \(\geqslant 1\) . The following conditions are equivalent:

\begin{enumerate}
\def\labelenumi{(\alph{enumi})}
\item
  The ideal of \(\mathbf{A}\) generated by the family \((x_{\lambda})\) is equal to \(\mathfrak{m}\) .
\item
  The family \((x_{\lambda})\) is a system of generators of the \(\mathbf{A}_0\) -algebra \(\mathbf{A}\) .
\item
  For all \(i \geqslant 0\) , the \(\mathbf{A}_0\) -module \(\mathbf{A}_i\) is generated by the elements of the form \(\prod_{\lambda} x_{\lambda}^{n_{\lambda}}\) which are of degree \(i\) in \(\mathbf{A}\) .
\end{enumerate}

Clearly conditions (b) and (c) are equivalent. If they hold, every element of m is of the form \(f((x_{\lambda}))\) where f is a polynomial of \(A_{0}[X_{\lambda}]_{\lambda \in L}\) with no constant term; then \(m = \sum_{\lambda \in L} Ax_{\lambda}\) , which proves that (c) implies (a). Conversely, suppose that condition (a) holds. Let \(A' = A_{0}[x_{\lambda}]_{\lambda \in L}\) be the sub- \(A_{0}\) -algebra of A generated by the family \((x_{\lambda})\) and let us show that \(A' = A\) . For this, it is sufficient to show that \(A_{i} \subset A'\) for all \(i \geqslant 0\) . We proceed by induction on i, the property being obvious for i = 0. Then let \(y \in A_{i}\) with \(i \geqslant 1\) . Since \(y \in m\) , there

exists a family \((a_{\lambda})_{\lambda \in \mathbf{L}}\) of elements of \(\mathbf{A}\) of finite support such that \(y = \sum_{\lambda} a_{\lambda} x_{\lambda}\) and we may assume that each of the \(a_{\lambda}\) is homogeneous of degree \(i - \deg(x_{\lambda})\) (by replacing it if need be by its homogeneous component of that degree); as \(\deg(x_{\lambda}) > 0\) , the induction hypothesis shows that \(a_{\lambda} \in \mathbf{A}'\) for all \(\lambda \in \mathbf{L}\) , whence \(y \in \mathbf{A}'\) and \(\mathbf{A}_t \subset \mathbf{A}'\) , which proves that (a) implies (b).

COROLLARY. Let \(\mathbf{A} = \bigoplus_{i\in \mathbb{Z}}\mathbf{A}_i\) be a graded commutative ring with positive degrees and m the graded ideal \(\bigoplus_{i\geqslant 1}\mathbf{A}_i\) .

\begin{enumerate}
\def\labelenumi{(\roman{enumi})}
\tightlist
\item
  The following conditions are equivalent:
\end{enumerate}

\begin{enumerate}
\def\labelenumi{(\alph{enumi})}
\item
  The ideal \(\mathfrak{m}\) is a finitely generated A-module.
\item
  The ring A is a finitely generated \(A_{0}\) -module.
\end{enumerate}

\begin{enumerate}
\def\labelenumi{(\roman{enumi})}
\setcounter{enumi}{1}
\item
  Suppose that the conditions in (i) hold and let \(\mathbf{M} = \bigoplus_{i\in \mathbf{Z}}\mathbf{M}_i\) be a finitely generated graded A-module. Then, for all \(i\in \mathbf{Z}\) , \(\mathbf{M}_i\) is a finitely \(\mathbf{A}_0\) -module and there exists \(i_0\) such that \(\mathbf{M}_i = \{0\}\) for \(i < i_0\) .
\item
  If a family \((y_{\mu})\) of elements of A is a system of generators of the A-module m (resp. of the \(A_{0}\) -module A), so is the family consisting of the homogeneous components of the \(y_{\mu}\) ; the equivalence of conditions (a) and (b) then follows from Proposition 1.
\item
  We may suppose that A is generated (as an \(\mathbf{A}_0\) -algebra) by homogeneous elements \(a_i\) ( \(1 \leqslant i \leqslant r\) ) of degree \(\geqslant 1\) and M is generated (as an A-module) by homogeneous elements \(x_j\) ( \(1 \leqslant j \leqslant s\) ); let \(h_i = \deg(a_i)\) , \(k_j = \deg(x_j)\) . Clearly \(\mathbf{M}_n\) consists of the linear combinations with coefficients in \(\mathbf{A}_0\) of the elements \(a_1^{\alpha_1} a_2^{\alpha_2} \ldots a_r^{\alpha_r} x_j\) such that the \(\alpha_i\) are integers \(\geqslant 0\) satisfying the relation \(k_j + \sum_{i=1} \alpha_i h_i = n\) ; for each \(n\) there is only a finite number of families \((\alpha_i)_{1 \leqslant i \leqslant r}\) satisfying these conditions, since \(h_i \geqslant 1\) for all \(i\) ; we conclude that \(\mathbf{M}_n\) is a finitely generated \(\mathbf{A}_0\) -module and moreover clearly \(\mathbf{M}_n = \{0\}\) when \(n < \inf_{j}(k_j)\) .
\end{enumerate}

